\documentclass[10pt,a4paper]{amsart}
\usepackage[margin=19mm]{geometry}
\usepackage{amsmath,amssymb,mathtools}
\usepackage[scr=rsfs,cal=euler]{mathalfa}
\usepackage{xfrac}
\usepackage[unicode,hidelinks,bookmarks=false]{hyperref}
\newcommand{\ie}{i.e.\ }
\newcommand{\cf}{cf.\ }
\newcommand{\resp}{respectively\ }
\newcommand{\Subsection}{\subsection}
\providecommand{\nobreakdash}{\nobreak\hbox{-}}
\setlength{\emergencystretch}{2em}
\multlinegap0pt
\usepackage{bm}
\usepackage{amssymb}
\usepackage[shortlabels]{enumitem}
\usepackage{float}
\usepackage{mathtools}
\newtheorem{lemma}{Lemma}
\newtheorem{theorem}{Theorem}
\title{Varieties with two smooth blow up structures}
\author{Supravat Sarkar}
\date{Source: arXiv:2409.10560v2 (CC BY 4.0)}
\begin{document}
\maketitle

\noindent\emph{Source and licence.} Varieties with two smooth blow up structures. Version arXiv:2409.10560v2 (CC BY 4.0), distributed by arXiv under the Creative Commons Attribution 4.0 International licence, http://creativecommons.org/licenses/by/4.0/, as recorded in the arXiv licence metadata for this exact version and shown on the abstract page https://arxiv.org/abs/2409.10560v2. Full licence notice: \texttt{licenses/arxiv-2409.10560-CC-BY-4.0.txt}.\par\smallskip

\noindent\textit{Editorial scope.} Counted unit: the article's theorem theorem:2case (source0.tex lines 263--269). The article's complete original proof of it is delivered with it (source0.tex lines 270--297, one \texttt{proof} environment of 28 lines). No mathematics is written, repaired or re-derived: the statement and the proof are the article's own lines, in the article's own notation, and no retained line is substituted or re-flowed.  Retained uncounted as the source-local context the proof needs: lines 168--176; lines 214--221; lines 248--250.  Nothing was omitted from any retained span, so the package carries no omission record.  No retained line contains a citation, so the wrapper carries no bibliography and no bibliography entry.  No retained line contains a figure environment, an image-inclusion command or drawing code, so no image is delivered and the package declares no asset dependency.  Editorial formatting only, in addition: an AMS \texttt{amsart} preamble that restates the article's own declarations and macros used by the retained lines, namely the environments \texttt{lemma}, \texttt{theorem} on separate counters, together with the standard packages those lines need.  The retained environments print in this document as Lemma 1, Theorem 1 in order of appearance, whereas the article prints the same items with the numbering of its own full document.  The complete native source of the article as distributed in the arXiv e-print for this exact version is bundled unchanged as \texttt{sources/165355/source0.tex}.
\begin{lemma}\label{lemma:katz}
    \begin{enumerate}
        \item [(i)] $(d-1)(n+1)=(n-m_1-1)\frac{cd-1}{a}$,
        \item [(ii)] $(c-1)(n+1)=(n-m_2-1)\frac{cd-1}{a}$,
        \item [(iii)] $\frac{c}{a}+\frac{1}{a}. \frac{n-m_2-1}{n-m_1-1}=\frac{n+1}{n-m_1-1}$,
        \item [(iv)]$\frac{d}{a}+\frac{1}{a} .\frac{n-m_1-1}{n-m_2-1}=\frac{n+1}{n-m_2-1}$,
        \item [(v)]$c>d\geq 2$.
    \end{enumerate}
\end{lemma}

\begin{lemma}\label{lemma:betti}
    \begin{enumerate}
        \item [(i)] $a_i\neq 0 \Longleftrightarrow 0\leq i\leq m_1$, $b_i\neq 0 \Longleftrightarrow 0\leq i\leq m_2$,
        \item[(ii)] $a_i-a_{i-(n-m_1-1)}=b_i-b_{i-(n-m_1-1)}$,
        \item[(iii)] $a_i=b_i$ if $0\leq i\leq n-m_1-2$,
        \item[(iv)] If $m_1\geq \frac{3n-2}{4}$, then $m_2\leq n-m_1-2$. 
    \end{enumerate}
\end{lemma}

\begin{theorem}\label{theorem:a=1}
In the notations of the setup, we have $a=1.$
\end{theorem}

\begin{theorem}\label{theorem:2case}
In the notations of the setup, we have
    \begin{enumerate}
        \item [either (1)] $n=4, a=1, c=3, d=2, m_1=2, m_2=1$,
     \item [or (2)] $n=9, a=1, c=3, d=2, m_1=6, m_2=4$.
    \end{enumerate}
\end{theorem}

\begin{proof}
    By Theorem \ref{theorem:a=1}, $a=1$. If $m_1\geq\frac{3n-2}{4}$, by Lemma \ref{lemma:betti} $(iv)$, $m_1+m_2\leq n-2$, and we have $$2\leq d<d+\frac{n-m_1-1}{n-m_2-1}=\frac{n+1}{n-m_2-1}\text{ (by Lemma \ref{lemma:katz}} (iv))$$ $$\leq\frac{n+1}{m_1+1}\leq\frac{n+1}{\frac{3n-2}{4}+1}=\frac{4n+4}{3n+2}.$$
    So, $$4n+4>2(3n+2)=6n+4\implies n<0,$$ a contradiction. So, $m_2<m_1<\frac{3n-2}{4}$.

    By Lemma \ref{lemma:katz} $(iii), (iv)$,
    \begin{equation}
        c=\frac{m_2+2}{n-m_1-1}, d=\frac{m_1+2}{n-m_2-1}.
    \end{equation}
    \textbf{Claim:} $c=3$, $d=2$.
    \begin{proof}
        By Lemma \ref{lemma:katz} $(v)$, it suffices to show $c\leq 3$. Suppose $c\geq 4$. By $(3)$, $m_2+2\geq 4(n-m_1-1)$, so $4n-6\leq m_2+4m_1<\frac{3n-2}{4}-1+3n-3$ (as $m_2\leq m_1-1<\frac{3n-2}{4}-1, 4m_1<3n-2\implies 4m_1\leq 3n-3$) $=15n/4-9/2$.

        Hence, $4n-6<15n/4-9/2$, that is, $n\leq 5$.

        If $n=5$, then $(m_2,m_1)=(1,2), (1,3)$ or $(2,3)$. In each case either $c$ or $d$ is not an integer, a contradiction.

        If $n=4$, we must have $m_1=2, m_2=1$. So, $c=3$ by $(3)$. Contradiction to $c\geq 4$.

        This contradiction shows $c\leq 3$.
    \end{proof}

    Now $(3)$ and claim gives 
    \begin{equation}
        m_1=\frac{4n-6}{5}, m_2=\frac{3n-7}{5}.
    \end{equation}
    Since $m_1, m_2$ are integers, we get $n\equiv -1(\text{mod } 5)$. Also, $$m_1<\frac{3n-2}{4}\implies \frac{4n-6}{5}<\frac{3n-2}{4}\implies n<14.$$ Together with $n\equiv -1(\text{mod } 5)$, we get $n=4$ or $9$.
    If $n=4$, then $m_1=2, m_2=1$ by $(4)$, so case $1)$ occurs. If $n=9$, then $m_1=6, m_2=4$ by $(4)$, so case $2)$ occurs. 
\end{proof}

\end{document}
